\summarychapter{8}{Apostles to Intellectuals}{The Alexandrians}

\begin{inbrief}
As Christianity grew into the chief rival of Rome's old ways, it had to decide what to do with Greek
philosophy and literature. Tertullian wanted nothing to do with them; at Alexandria Clement and Origen
showed that faith and philosophy could live together when both bow to Christ, and so prepared the way
for a Christian empire and for later theology.
\question{Must Christians resist Greek philosophy as an enemy, or can it be made an ally of the gospel?}
\end{inbrief}

About 374 the young monk Jerome, later the translator of the Latin Vulgate, dreamed during an illness
that he stood before the judgment seat and heard a voice: ``You are a Ciceronian, not a Christian.''
He loved Christ, but he also loved Cicero, Virgil and Horace. His nightmare is typical of the early
church's struggle over pagan literature and philosophy, a struggle that has never ended. As the faith
spread into the world of Hellenism, new arguments and criticisms pushed Christian thinkers toward a
fuller account of their teaching, and people of culture and power were asking what part Christianity
would play in the affairs of men and empires.

\ssection{Light in the World}
The church always stands in a double relation to the world: ``not of the world'' yet ``sent into the
world.'' It moves to a rhythm of detachment, because the gospel comes from God, and involvement,
because God sends it to be light. Christians differ over where the boundary lies; what is witness to
one is compromise to another. Some therefore resisted every attempt to reconcile the faith with Greek
philosophy, appealing to John's warning not to love the world and to Paul's saying that the cross is
foolishness to Greeks. The fiercest was Tertullian. Philosophy, he cried, breeds heresy: Valentinus
was a Platonist, Marcion a Stoic. ``What do Athens and Jerusalem have in common?'' Once we believe,
we need nothing further. In view of the life-and-death struggle with Gnosticism, a Christianity recast
to fit fashionable ideas, this was no mere bigotry.

Yet during the very century of persecution Christians were finding ways to state the gospel in
Hellenic terms, and those teachers paved the way for Rome's later acceptance of the faith. The lead
came from the catechetical school of Alexandria under Clement and Origen, the first of a line of
scholars at home in Greek wisdom yet loyal to Christ, who tried to bring into Christianity the best of
Hellenic culture, especially of the Platonists and Stoics. ``The way of truth is one,'' said Clement,
``but into it, as into a perennial river, streams flow from all sides.'' Shelley thinks the school
began independently of the church, like the schools philosophers opened in the great cities and like
Justin's lodging in Rome, where he taught all who came. About 180 the Sicilian Pantaenus founded such
a school of ``Christian Gnosticism'' in Alexandria: gnostic because it asked the great questions of
meaning, Christian because it kept orthodox answers. He won many pagans and deepened many Christians.

\ssection{A Pastor to Philosophers}
Pantaenus drew Clement to Alexandria, where he stayed twenty years as pupil and then as head of the
school. His chief surviving works form a kind of trilogy: the \emph{Exhortation to the Heathen}, the
\emph{Instructor} and the unfinished \emph{Miscellanies}. His contemporaries knew him as a messenger of
Christianity in a philosopher's cloak. ``The first Christian scholar,'' he knew the Scriptures and the
learning of his time, and he understood the students who came from Rome, Athens and Antioch
dissatisfied with their teachers, many of them already touched by some gnostic theory. He took up not
only the dress and language of the philosophers but their problems --- the origin of the world, evil,
salvation through the Word --- and showed that these find their deepest answer in the Christian
revelation. His aim was pastoral, not merely theological: to win people, not arguments.

It was a risky venture in a city full of Valentinian Gnosticism, where converts often seemed to face a
choice between a clever heresy and a modest orthodoxy. Clement offered a third way. He gave philosophy
a positive role as a preparation for Christ: it was ``a schoolmaster to the Greeks, as the law was to
the Hebrews.'' But his method was the opposite of the heretics'. They wrecked the apostolic gospel; he
and Origen stayed loyal to the message of Peter and Paul. They also differed on conduct. Many Gnostics
separated knowledge from character; Clement taught that insight comes to the pure in heart, who love
the good for its own sake, in an ascent from faith through knowledge to the vision of God beyond this
life. The decisive difference was creation: the Valentinians rejected it as the work of an evil deity,
while for Clement God planted seeds of truth in all rational creatures, so that Christians can learn
from the Greeks, since all truth and goodness come from the Creator. After Clement, Greek and
Christian thought were joined; without that union the theology of the first great councils would have
been impossible.

\ssection{Origen and the Thirst for Truth}
When persecution struck Alexandria shortly after 200, Clement fled and left the school to a gifted
youth of eighteen. Origen came from a Christian home; his father Leonides died a martyr in that
persecution, and Origen sold his secular books to support the family. He shared Clement's love of
philosophy and held that the desire to know the causes of things was planted in us by God. He was
often at odds with his bishop, Demetrius, whom he thought power-hungry and who found him
uncooperative. About 229, on his way to Athens, Origen was ordained at Caesarea in Palestine;
Demetrius treated this as insubordination and had him publicly condemned. So his career falls into
the years at Alexandria (202--230) and at Caesarea (230--254).

He was a magnetic teacher to whom students traveled hundreds of miles, among them Gregory, later
called the Wonder-worker for his missionary success in Asia Minor, who wrote a book in his praise.
Origen set before his pupils the true goal of philosophy, the good life; ignorance, he said, is a
barrier to godliness, and true philosophy centers on the Word. He formed character as well as ideas
and taught, Gregory says, more by his acts than by his theories. A wealthy friend paid for seven
shorthand writers, and books poured out; Jerome asked who could ever read them all. He regarded the
exposition of Scripture as his first task. Almost alone among Greek-speaking Christians he studied the
Old Testament in Hebrew. His \emph{Hexapla} set six versions of it in parallel columns, and he wrote
scores of commentaries and hundreds of sermons. Where the surface sense of a passage seemed unworthy
of God, he took it as a sign of a deeper lesson. Hence his ``allegorical interpretation'' with three
levels of meaning, literal, moral and spiritual, matching body, soul and mind. Above all he wanted
the whole Bible to speak, against the heretics' habit of forcing a few striking texts. This work let
intelligent Christians go on believing the Bible. Origen, Shelley says, saved the Scriptures for the
church and protected the historical foundation of the faith.

\ssection{A Theology for Thinkers}
Origen was also the first to set out a whole intellectual framework of the faith rather than merely
refute heresy. His \emph{First Principles}, an early work he never needed to revise much, was written
for educated readers: he did not despise the faith of simple people, but he saw that Christianity
must justify itself to the mind as well as the heart if it was to shape civilization. He held to the
rule of faith and felt free to speculate where it was silent. His vision reached so far as to teach
that all creatures, even the devil, would at last be restored to God and hell emptied. This doctrine
brought him endless trouble. His error, Shelley judges, was to turn a dream into a doctrine: if God
has character and man has choice (as Charles Williams argued), an everlasting rejection of God must
remain possible.

In the persecution of Decius Origen was singled out, imprisoned, chained and tortured, but kept alive.
Released when the terror ended in 251, he died of its effects at Tyre in 254, aged sixty-nine. Did
Clement and Origen go too far and let Greek ideas take the gospel captive? Shelley answers that they
were true to the faith as it was then understood: as Christ took on human existence, so his people
take on the humanity of every civilization. ``For the Hellenes one must become a Hellene,'' said
Clement. They preserved humanism for Christianity, made possible the later work of Athanasius, Gregory
of Nyssa and John Chrysostom, and showed that the best of classical culture had a home and a future in
the church.

\begin{keyfacts}
\keypeople{Jerome; Tertullian; Valentinus; Justin; Pantaenus; Clement of Alexandria; Origen; Leonides;
Demetrius, bishop of Alexandria; Gregory the Wonder-worker; Decius.}
\keydates{%
  \kd{c. 180}{Pantaenus opens his Christian school at Alexandria}
  \kd{185}{Birth of Origen}
  \kd{c. 202}{Persecution at Alexandria; Clement flees and Origen, aged 18, heads the school}
  \kd{c. 229--230}{Origen ordained at Caesarea, condemned by Demetrius, settles in Palestine}
  \kd{254}{Death of Origen at Tyre after torture under Decius}
  \kd{c. 374}{Jerome's dream: ``a Ciceronian, not a Christian''}}
\keyterms{Hellenism; catechetical school of Alexandria; Christian Gnosticism; philosophy as
schoolmaster; \emph{Hexapla}; allegorical interpretation (literal, moral, spiritual); \emph{First
Principles}; universal restoration.}
\end{keyfacts}
